\clearpage
\pagestyle{empty}
\chapter*{{\emph\LARGE{General Introduction}}}
%\chapter*{General introduction}
%\chaptermark{General introduction}

\markboth{General introduction}{}
\addcontentsline{toc}{chapter}{General introduction}

\pagestyle{fancy}
\fancyhf{}
\fancyhead[HL]{General introduction}
\cfoot{\thepage}

\lettrine[lines=2,loversize=0.1]{\textbf{\textit{I}}} 
ndustrial systems are increasingly relying on connected devices capable of acquiring, processing, and exchanging information in real time. In such environments, sensors, actuators, and supervisory systems must interact efficiently through various communication interfaces and protocols. However, integrating these heterogeneous components often requires significant development effort, particularly when flexibility, scalability, and interoperability are required.

To address these challenges, Be Wireless Solutions (BWS) launched the development of the IOBox platform, a configurable industrial Input/Output system designed to serve as a universal interface between field devices and monitoring or control systems. The platform is intended to simplify the integration of sensors and actuators while providing data acquisition, signal processing, control, and communication functionalities within a single solution.

The IOBox supports multiple industrial communication protocols and offers both analog and digital input/output capabilities. Its modular architecture enables adaptation to various industrial applications, ranging from environmental monitoring and energy management to process automation and equipment supervision. By centralizing these functionalities in a single platform, the IOBox reduces deployment complexity and accelerates the development of industrial IoT solutions.

The objective of this end-of-study project is to participate in the design and implementation of the IOBox platform. The work focuses on the development of embedded software components, communication interfaces, peripheral drivers, sensor integration, and system architecture required to ensure reliable operation in industrial environments.\linebreak[4]

This report, which crowns the work entrusted to us, is structured around three chapters:

\textbf{Chapter 1} establishes the general context of the project. It begins by presenting the academic institution and the host company, Be Wireless Solutions, describing its organization, its main activities, and the industrial IoT solutions it provides. It then defines the problem statement that motivated this work, namely the need to replace the custom firmware developed for each new setup with a single reusable base, and introduces the IOBOX platform proposed as a solution. The chapter closes by stating the project objectives and detailing the V-Model methodology adopted to structure the work from requirements analysis through to system testing.

    \textbf{Chapter 2} introduces the general concept of the IOBOX platform and presents the technical environment on which it is built. It first describes the overall hardware architecture and its central component, the STM32G474 microcontroller, followed by the external EEPROM and the sensors integrated into the system. It then presents the on-chip peripherals used for acquisition and control, including the ADC, DAC, timers, and PWM, before detailing the communication protocols supported by the platform, such as I2C, SPI, UART, RS485, Modbus, CAN, and LIN. Finally, it introduces the software tools used throughout the development of the project.

    \textbf{Chapter 3} details the design and practical implementation of the platform. It presents the layered software architecture, organized into the HAL, BSP, Manager, and Application layers, and explains the role of each within the system. It then describes the development of the board support package and peripheral drivers, the management of the communication interfaces, the integration and configuration of the sensors, and the construction of the output data frame. The chapter concludes with the functional testing performed on the target hardware, where each peripheral, sensor, and communication interface is tested individually to confirm its correct behavior.
